\documentclass[a4paper,12pt]{article}
\usepackage{codeanatomy}
\usepackage{pgfplotstable}
\pgfplotsset{compat=1.18}
\usepackage{multicol}
\usepackage{booktabs}
\usepackage{amsmath}
\usepackage{graphicx}
\usepackage{float}
\usepackage{listings}
\usepackage{mathtools}
\usepackage[thinc]{esdiff}
\usepackage[T2A]{fontenc} 
\usepackage[utf8]{inputenc}
\usepackage{titlesec, titletoc} 
\usepackage{cmap} 
\usepackage[english,russian]{babel} 
\usepackage{caption}
\usepackage{xcolor} 
\usepackage{lastpage}
\usepackage{geometry}
    \geometry{left=30mm, right=20mm, top=20mm, bottom=20mm}

\titlespacing*{\section}{0pt}{2\baselineskip}{\baselineskip}
\titlespacing*{\subsection}{0pt}{2\baselineskip}{\baselineskip}
\titlespacing*{\subsubsection}{0pt}{2\baselineskip}{\baselineskip}


\definecolor{codegreen}{rgb}{0,0.6,0}
\definecolor{codegray}{rgb}{0.5,0.5,0.5}
\definecolor{codepurple}{rgb}{0.58,0,0.82}
\definecolor{backcolour}{rgb}{0.95,0.95,0.92}

\lstdefinestyle{mystyle}{
    backgroundcolor=\color{backcolour},   
    commentstyle=\color{codegreen},
    keywordstyle=\color{magenta},
    numberstyle=\tiny\color{codegray},
    stringstyle=\color{codepurple},
    basicstyle=\ttfamily\footnotesize,
    breakatwhitespace=false,         
    breaklines=true,                 
    captionpos=b,                    
    keepspaces=true,                 
    numbers=left,                    
    numbersep=5pt,                  
    showspaces=false,                
    showstringspaces=false,
    showtabs=false,                  
    tabsize=2
}

\lstset{style=mystyle}

% ===================== ПЕРЕМЕННЫЕ =====================
\newcommand{\University}{Санкт-Петербургский политехнический университет Петра Великого}
\newcommand{\Institute}{Институт компьютерных наук и кибербезопасности}
\newcommand{\School}{Высшая школа программной инженерии}

\newcommand{\WorkType}{КУРСОВОЙ ПРОЕКТ}
\newcommand{\Topic}{Моделирование динамической сцены}
\newcommand{\Variant}{33} 
\newcommand{\Discipline}{Компьютерная графика}

\newcommand{\Group}{в5130904/30322}
\newcommand{\StudentName}{Соловей Г. И.}
\newcommand{\SupervisorName}{Леонтьева Т. В.}

\newcommand{\City}{Санкт-Петербург}
\newcommand{\Year}{2026}

% --- Задание ---
\newcommand{\Deadline}{20.04.2026}
\newcommand{\StartDate}{02.02.2026}
\newcommand{\Software}{Visual Studio Code}
\newcommand{\Libraries}{OpenGL}

% =====================================================

\begin{document}

\begin{center}
    {\small \University}\\
    {\small \Institute}\\
    {\small \School}

    \vspace{4cm}

    {\large \textbf{\WorkType}}\\[1cm]

    {\textbf{\Topic, вариант №\Variant}}\\

    по дисциплине «\Discipline»

    \vspace{3cm}
\end{center}

\begin{flushleft}
    Выполнил\\
    студент гр. \Group \hfill \StudentName

    \vspace{1cm}

    Руководитель \hfill \SupervisorName
\end{flushleft}

\vfill

\begin{center}
    \City\\
    \Year
\end{center}

\thispagestyle{empty}

\newpage

\begin{center}
    \textbf{ЗАДАНИЕ \\ НА ВЫПОЛНЕНИЕ КУРСОВОГО ПРОЕКТА}

\end{center}
студенту группы \Group \hfill \StudentName

\vspace{0.5cm}

\begin{enumerate}
    \item \textbf{Тема проекта (работы):} \Topic, вариант №\Variant
    \item \textbf{Срок сдачи студентом законченного проекта (работы):} \hfill 20.04.2026
    \item \textbf{Исходные данные к проекту (работе):} Задание с параметрами согласно варианту, среда разработки (\Software),
          графические библиотеки (\Libraries)
    \item \textbf{Содержание пояснительной записки:}
          задание,
          описание алгоритма,
          описание реализации,
          скриншоты,
          список литературы,
          приложение (код программы) \\
          Объем пояснительной записки: \pageref{LastPage} страниц
    \item Перечень графического материала ( с указанием обязательных чертежей и
          плакатов)	------
    \item Консультанты ------
    \item Дата получения задания: \StartDate
\end{enumerate}

\vspace{1cm}

\noindent
Руководитель \hfill \SupervisorName

\vspace{1cm}

\noindent
Задание принял к исполнению \hfill \StudentName

\vspace{1cm}

\hfill \StartDate

\thispagestyle{empty}

\newpage

\tableofcontents
\setcounter{page}{1}
\thispagestyle{empty}

\newpage

\section{Описание задания}
\subsection{Описание области задачи}

Системы частиц (particle system) - это объекты-эмиттеры (emitters), генерирующие по
заданному алгоритму частицы точечной массы с определенной формой, начальной
скоростью, сроком жизни и другими свойствами.
Такие анимационные объекты используются для моделирования дождя, огня, дыма,
пузырьков газа в жидкости, осколков взрывающихся снарядов и тому подобных образцов
объективной реальности. С помощью систем частиц также составляются модели
поведения групп живых существ - стай птиц, косяков рыб.
Время в системе частиц дискретно. За один такт времени каждая из частиц может
переместиться под влиянием других частиц или внешней среды (т.е. каких-то сил,
определенных не для каждой частицы отдельно, а для всей системы в целом - например,
силы гравитации). Поведение системы задается правилами поведения частиц в данной
системе.
Количество частиц в системе может быть постоянным либо меняться со временем.
Если количество частиц может меняться, то могут появляться (зарождаться) новые
частицы или исчезать (умирать) старые.
Свойствами частицы могут быть, например, цвет, время жизни, возраст, размер,
скорость и направление движения и т.д. Набор свойств определяет тип частицы. Эти
свойства могут меняться в течение жизни частицы, но сам набор (состав) свойств частицы
остается неизменным. При зарождении частицы ее свойства инициализируются
начальными значениями.

\subsection{Описание общей части задания}

\begin{enumerate}
    \item  В качестве источников-генераторов (еmitters) частиц выступают объекты различной
          конфигурации. Количество частиц от 100 до 300. Координаты начальной позиции частицы
          рассчитываются при помощи случайного числа, переведенного на поверхность требуемой
          фигуры. Первоначально направление движения частицы определяется по нормали к
          поверхности фигуры.
    \item Каждая частица должна характеризоваться следующим набором параметров:
          координаты, скорость, размер, цвет (с возможным добавлением текстуры и прозрачности),
          время жизни, след от частицы. Параметры могут зависеть друг от друга (например, цвет
          частицы изменяется в зависимости от времени жизни). По окончании времени жизни
          частица исчезает и порождается новая частица. Тем самым система становится
          цикличной.
    \item Так как время системы дискретно, то определяются правила поведения частиц на
          каждом такте (dtime). Частицы должны изменять свое расположение по определенным
          законам, одинаковым для всей системы. Новая позиция частицы рассчитывается по
          следующим правилам (в векторной форме):
          \begin{equation}
              \vec{v} = \vec{v} + \vec{a} * dtime
          \end{equation}
          \begin{equation}
              \vec{pos} = \vec{pos} + \vec{v} * dtime
          \end{equation}
          Возможно использование собственных идей для расчета новой позиции частицы.
          Поощряется красивая формула для расчета изменений.
    \item За каждой из частиц остается след, который изображается при помощи повторения
          частицы (поощряются также другие способы изображения следа). След может иметь
          различную протяженность.
    \item Требуется осуществить моделирование столкновений или/и физических сил,
          воздействующих на систему.
          Аттракторы: Для моделирования физических воздействий используются
          «магниты»-притягиватели (аттракторы) в форме стандартных примитивов (Cylinder, Point,
          Plane). При приближении к таким объектам частицы должны изменять скорость и
          направление движения в зависимости от расстояния до него. В случае с аттрактором:
          когда расстояние до аттрактора становится меньше заданного, частица изменяет
          направление движения и притягивается к аттрактору (по принципу сложения векторов, с
          учетом нормали к объекту-аттрактору) с сохранением или увеличением скорости.
\end{enumerate}

\subsection{Индивидуальное задание}
\subsection*{Задание 33}
\begin{enumerate}
    \item Эмиттер – тетраэдр
    \item Обязательные параметры: размер увеличивается с увеличением расстояния от эмиттера.
          Остальные параметры устанавливаются и изменяются по вашему выбору.
    \item След: присутствует, длина от 4 до 10
    \item Аттрактор: плоскость
\end{enumerate}

\newpage

\section{Выполнение проекта}
\subsection{Выбор средств разработки}

Для выполнения задания были выбраны следующие библиотеки и инструменты:
\begin{itemize}
    \item библиотека OpenGL v.4.2 для C++;
    \item библиотека GLFW для создания и управления окном;
    \item библиотека GLM для работы с линейной алгеброй;
    \item библиотека с открытм исходным кодом ImGui для отрисовки простых элементов интерфейса;
    \item пакетный менеджер Conan 2.0 для управления зависимостями;
    \item система сборки CMake.
\end{itemize}
Библиотеку GLUT (FreeGLUT) было решено не использовать, так как она предоставляет функции устаревшего на
момент выполнения работ и проекта пайплайна графической отрисовки (fixed-function pipeline). В современном OpenGL большая часть рендеринга
возлагается на видеокарту посредством шейдеров, чтобы минимизировать использование процессора в пайплайне отрисовки.

\subsection{Библиотека-движок}
Для выполнения всех лабораторных работ по дисциплине была разработана отдельная графическая библиотека на основе используемых инструментов.
Разработанная библиотека представляет собой упрощённый графический движок,
реализующий базовые компоненты современного графического приложения: управление окном,
камерой, сценой, её объектами и процессом рендеринга.

В данном подразделе будет представлено краткое описание структуры библиотеки, листинг исходного кода представлен в приложении А.

\subsection*{Описание структуры}
\subsubsection*{Application}
Application — центральный синглтон-класс движка, отвечающий за жизненный цикл приложения. Управляет созданием окна (GLFW), инициализацией графики, запуском основного цикла и хранением текущей сцены. Выступает точкой входа и координирует взаимодействие между системой рендера и сценами.

\subsubsection*{Scene}
Scene — контейнер графической сцены, управляющий набором объектов и активной камерой. Отвечает за обновление объектов (tick) и их жизненный цикл (создание/удаление). Используется Application как текущая сцена для обработки и рендеринга.

\subsubsection*{Camera}
Camera — класс, описывающий параметры наблюдения сцены. Хранит позицию, направление взгляда и параметры проекции (perspective/orthographic), а также вычисляет матрицы вида и проекции. Используется для определения того, как сцена отображается на экране.

\subsubsection*{Object/RenderObject}
Object / RenderObject -- базовая иерархия сущностей сцены.
Object задаёт общий интерфейс объектов (обновление, удаление, клонирование) и связь со сценой.
RenderObject расширяет его, добавляя трансформацию и интерфейс рендеринга, представляя объекты, участвующие в отрисовке.

\subsubsection*{Mesh/OBJMesh}
Mesh / OBJMesh — реализация рендерируемых объектов геометрии.
Mesh отвечает за отрисовку, хранит геометрию, GPU-буферы и материал.
OBJMesh — конкретная реализация, загружающая геометрию из .obj файла и подготавливающая её для рендеринга.

\subsubsection*{Material}
Material — класс, отвечающий за шейдеры и их использование при рендеринге. Хранит исходный код вершинного и фрагментного шейдеров, компилирует их в GPU-программу, проверяет ошибки компиляции и предоставляет метод use() для активации шейдера при отрисовке объектов.

\subsection{Выполнение задания курсового проекта}
% \begin{figure}[H]
%     \centering
%     \includegraphics[width=\textwidth]{images/image.png}
%     \caption{Первичная сцена приложения}
% \end{figure}
% \subsubsection*{textured\_material.cc}
% \lstinputlisting{../src/c++/colored_material.cc}
% \newpage
\end{document}
